% ============================================================================
% Hypothesis H3: robust versus nominal site selection
% Pre-registered decision rule (OQ-050, H3), fixed 2026-10-13, before any
% real-country result was produced. Primary threshold delta = 0.20, m = k.
% Notation follows docs/METHODOLOGY.md and the H1 and H2 sections.
% Citation keys marked "verify" must be checked before submission.
% ============================================================================

\section{Hypothesis H3: Robust Versus Nominal Site Selection}
\label{sec:h3}

\subsection{Purpose and Position in the Research Design}
\label{subsec:h3_purpose}

Hypothesis H3 asks whether planning for robustness changes where renewable
capacity would be built. It compares two shortlists of the same size: the sites
that are best under present climate and nominal costs, which a conventional
assessment would select, and the sites that are best when climate and
techno-economic uncertainty are considered together, measured by the maximum
relative regret. If the two shortlists largely coincide, a conventional
assessment already identifies robust sites and the added complexity of a
robustness analysis changes little for siting. If they differ materially, a plan
based on present conditions would place a substantial part of the capacity in
locations that are not robust.

The question is a spatial instance of a general contrast in decision analysis
under uncertainty: a strategy that is optimal for a single best-estimate future
need not perform well across the range of plausible futures, and a robust
strategy may accept some loss of performance in the nominal case in exchange
for lower sensitivity to assumptions
\cite{lempert2003,lempertcollins2007}. In optimisation, the loss accepted in the
nominal case is known as the price of robustness \cite{bertsimassim2004}; in
stochastic programming, the benefit of planning with the uncertainty rather than
with a single expected scenario is the value of the stochastic solution
\cite{birge1982}. In energy system modelling, the existence of many
near-optimal and structurally different solutions shows that a single optimum
is a weak basis for siting advice \cite{decarolis2011,neumannbrown2021}.
% verify decarolis2011

H3 is distinct from the two preceding hypotheses. H1b starts from the
present-optimal sites and asks how many are lost when only the climate changes;
H3 starts from the robust sites, under climate and techno-economic uncertainty
together, and asks how much of them a present-based plan would have missed. H2
compares two spatial scales under the same criterion; H3 compares two criteria
at the same scale.

\subsection{Notation}
\label{subsec:h3_notation}

The notation of Sections~\ref{sec:h1} and~\ref{sec:h2} applies. Both rankings
are computed over the same set $\mathcal{I}$ of cells that enter the robustness
analysis, so that any difference between them reflects the criterion and not the
definition of the candidate set. For each cell $i\in\mathcal{I}$, $q_i$ is the
installable capacity of the central land scenario and
$z_i=\mathrm{MR}_i$ is the primary robustness score.

The shortlist size is $k=\lceil p_k N\rceil$ with $N=|\mathcal{I}_0|$, as in H1
and H2. The tolerance is $m=k$. The sets used are:
\begin{itemize}
    \item $S_k$: the $k$ cells of $\mathcal{I}$ with the lowest $z_i$ (the
    robust shortlist, identical to the one that defines the target of H2);
    \item $r^{\mathrm{nom}}_i$: the rank of cell $i$ within $\mathcal{I}$ by
    $\mathrm{LCOE}_{i,f_0}$, the levelised cost under the reference climate and
    nominal parameters;
    \item $S^{\mathrm{nom}}_k$: the $k$ cells of $\mathcal{I}$ with the lowest
    $\mathrm{LCOE}_{i,f_0}$ (the nominal shortlist).
\end{itemize}
Ranks follow the deterministic tie convention of the F7 phase record, applied
identically to both criteria.

\subsection{Primary Statistic: Robust Capacity Outside the Nominal Plan}
\label{subsec:h3_statistic}

The statistic is the share of the robust capacity located in cells that a
present-based plan would not reach even with the tolerance:
\begin{equation}
    D=\frac{\sum_{i\in S_k}q_i\,\mathbb{1}\!\left[r^{\mathrm{nom}}_i>k+m\right]}
    {\sum_{i\in S_k}q_i}.
    \label{eq:h3_statistic}
\end{equation}
A robust cell ranked between $k+1$ and $k+m$ in the nominal ranking is not
counted: it lies just outside the nominal shortlist, at a small cost margin, and
its exclusion is boundary fluctuation rather than a different siting choice. The
tolerance $m=k$ is the same as in H1b. Capacity weighting follows H2, so that the
statistic measures where megawatts would be built rather than how many cells
change.

\subsubsection{Reference value}

If the robust and nominal rankings were unrelated, the nominal rank of a robust
cell would be uniform over $\mathcal{I}$, and the expected share would be
\begin{equation}
    D_{\mathrm{ref}}=1-\frac{k+m}{|\mathcal{I}|}.
    \label{eq:h3_reference}
\end{equation}
For $p_k=0.10$ this is close to $0.8$. Since the quality of the solar or wind
resource enters both criteria, the observed $D$ is expected to lie well below
$D_{\mathrm{ref}}$; how far below is the substance of the result.
$D_{\mathrm{ref}}$ is reported with every result.

\subsection{Hypothesis and Decision Rule}
\label{subsec:h3_rule}

H3 states that a material share of the robust capacity lies outside the reach of
a present-based plan. For each country and technology:
\begin{equation}
    \text{H3 is}\;
    \begin{cases}
        \text{not testable}, & D_{\mathrm{ref}}<\delta,\\
        \text{supported}, & D_{\mathrm{ref}}\geq\delta \text{ and } D\geq\delta,\\
        \text{not supported}, & D_{\mathrm{ref}}\geq\delta \text{ and } D<\delta,
    \end{cases}
    \qquad \delta=0.20 .
    \label{eq:h3_rule}
\end{equation}

The threshold $\delta=0.20$ reads as one in five robust megawatts placed where a
present-based plan would not go, and it is the same proportion that defines
vulnerability in H1b. The not-testable case covers shortlists so large that even
unrelated rankings could not reach $\delta$; for the top-$k$ shares used in the
thesis it does not occur, but it is checked and reported. As for H1 and H2, the
threshold is a decision of this thesis: the literature on robust and optimal
strategies motivates the comparison but provides no universal cut-off for a
material difference. The defence rests on its pre-registration, its direct
reading and its consistency with H1b.

\subsection{Magnitude: Price and Benefit of Robustness (No Verdict)}
\label{subsec:h3_magnitude}

Whether the two shortlists differ says nothing, by itself, about how much is at
stake. Two capacity-weighted quantities translate the difference into terms a
planner can use. For a set $S$ of cells, let
$\overline{X}(S)=\sum_{i\in S}q_iX_i\,/\sum_{i\in S}q_i$ denote the
capacity-weighted mean of a cell quantity $X_i$.

\subsubsection{Price of robustness}

\begin{equation}
    \Pi=\frac{\overline{\mathrm{LCOE}}_{f_0}(S_k)
    -\overline{\mathrm{LCOE}}_{f_0}(S^{\mathrm{nom}}_k)}
    {\overline{\mathrm{LCOE}}_{f_0}(S^{\mathrm{nom}}_k)}.
    \label{eq:h3_price}
\end{equation}
$\Pi$ is the relative increase of the present levelised cost incurred by
choosing the robust shortlist instead of the nominal one.

\subsubsection{Benefit of robustness}

\begin{equation}
    B=\overline{\mathrm{MR}}(S^{\mathrm{nom}}_k)-\overline{\mathrm{MR}}(S_k).
    \label{eq:h3_benefit}
\end{equation}
$B$ is the reduction of the worst-case relative regret obtained by choosing the
robust shortlist, in the units of $\mathrm{MR}$.

The pair $(\Pi,B)$ is the main descriptive result of H3. A small $\Pi$ with a
large $B$ means that robustness is cheap; a large $\Pi$ with a small $B$ means
that it is expensive for little protection. Both are conditional on the futures
and parameter ranges of the analysis, and $\Pi$ is evaluated in the nominal
future only. The capacities of the two shortlists, which may differ because cells
differ in capacity, are reported alongside.

\subsection{Complementary Diagnostics (No Verdict)}
\label{subsec:h3_diagnostics}

\begin{itemize}
    \item Location. The geodesic distance between the capacity-weighted centroids
    of $S_k$ and $S^{\mathrm{nom}}_k$; the total variation distance between their
    distributions of capacity across Admin1 units, defined as in H2; and a map of
    the two shortlists.
    \item Ranking. The Spearman rank correlation between $\mathrm{LCOE}_{i,f_0}$
    and $z_i$ over $\mathcal{I}$, and the Jaccard overlap of $S_k$ and
    $S^{\mathrm{nom}}_k$ by cell count.
\end{itemize}

The attribution of the difference to the climate or to the techno-economic axis
is the object of H4 and is not repeated here.

\subsection{Sensitivity Analyses}
\label{subsec:h3_sensitivity}

The verdict is issued from the primary configuration only; a change of verdict
in a variant is reported as a dependence of the result on that choice.
\begin{enumerate}
    \item Without IPSL-CM6A-LR, the climate model retained as a recorded
    exception.
    \item The 2071--2100 window.
    \item The 0.1-degree cell scale.
    \item Top-$k$ share $p_k\in\{0.05,0.10,0.20\}$.
    \item Tolerance $m\in\{k/2,\,2k\}$.
\end{enumerate}

\subsection{Worked Example}
\label{subsec:h3_example}

A country has $|\mathcal{I}|=1{,}000$ cells and $k=50$, so $m=50$ and
$D_{\mathrm{ref}}=1-100/1{,}000=0.90\geq0.20$; the test is informative. The
robust shortlist $S_k$ holds 500~MW. By nominal rank:
\begin{itemize}
    \item 300~MW lie in cells ranked 1 to 50 nominally, which a present-based
    plan would also select;
    \item 50~MW lie in cells ranked 51 to 100, just outside the nominal
    shortlist; they are not counted;
    \item 150~MW lie in cells ranked beyond 100, which a present-based plan
    would not select.
\end{itemize}
Then $D=150/500=0.30\geq0.20$ and H3 is supported. If, in addition, the
capacity-weighted present LCOE is 40.8 for the robust shortlist and 40.0 for the
nominal one, and the capacity-weighted maximum relative regret is 0.10 and 0.25
respectively, then $\Pi=0.02$ and $B=0.15$: building the robust shortlist costs
two per cent more under present conditions and reduces the worst-case regret by
fifteen percentage points. The values are illustrative and are not GeoFREA
results.

\subsection{Parameters}
\label{subsec:h3_parameters}

\begin{table}[htbp]
    \centering
    \caption{Pre-registered parameters of H3.}
    \label{tab:h3_parameters}
    \small
    \renewcommand{\arraystretch}{1.2}
    \begin{tabularx}{\textwidth}{
        @{}
        >{\raggedright\arraybackslash}p{2.4cm}
        >{\raggedright\arraybackslash}p{3.0cm}
        >{\raggedright\arraybackslash}X
        @{}
    }
        \toprule
        \textbf{Parameter} & \textbf{Value} & \textbf{Role} \\
        \midrule
        $\delta$ & 0.20 & Primary threshold; H3 is supported when at least one in
        five robust megawatts lies beyond the nominal top $k+m$. \\
        \addlinespace
        $m$ & $k$ & Tolerance beyond the nominal shortlist, as in H1b. \\
        \addlinespace
        $k$ & $\lceil p_k N\rceil$, $N=|\mathcal{I}_0|$ & Shortlist size shared
        with H1 and H2; $p_k$ fixed in OQ-021. \\
        \addlinespace
        Domain & $\mathcal{I}$ & Both rankings over the same candidate set. \\
        \addlinespace
        Weight & $q_i$ & Central-scenario capacity, as in H2. \\
        \addlinespace
        $D_{\mathrm{ref}}$ & $1-(k+m)/|\mathcal{I}|$ & Expected share under
        unrelated rankings; the test is not testable when
        $D_{\mathrm{ref}}<\delta$. \\
        \bottomrule
    \end{tabularx}
\end{table}

\subsection{Limitations}
\label{subsec:h3_limitations}

Shared drivers. The quality of the resource enters both the nominal cost and the
robustness score, so the two rankings are correlated by construction. H3 measures
how much of the difference survives this common driver; a not-supported verdict
is therefore an informative result, indicating that present-day assessment
already captures most of what robustness requires for siting.

Dependence on the robustness score. The robust shortlist is defined by
$\mathrm{MR}_i$, the maximum over climate members of the 90th percentile over
parameter samples. Another robustness metric could select a different shortlist;
the agreement between the regret and satisficing rankings is reported separately
as part of the robustness analysis.

Scope of the price of robustness. $\Pi$ is computed under present climate and
nominal parameters only. It is a statement about the present cost of a robust
choice, not about system costs, which would require system-level modelling.

Capacity differences. The two shortlists have the same number of cells but not
necessarily the same capacity; the statistic and the magnitudes are
capacity-weighted, and the two capacities are reported.

Generalisation. Results are conditional on the country, technology, candidate
set, futures and parameter ranges, and are reported per country and technology.
